\backmatter
\AfterBibliographyPreamble{%
Every measurement in this work is built on somebody's published result, and the
first purpose of this chapter is that those people are credited by name. A surname
in a footnote is not a citation. The second purpose is that the credit is checkable:
each source is registered in \texttt{orior\_citations}, where the copy sits
itself, and this bibliography is generated from that registry and pinned to one
commit of it. A reference therefore resolves not to a title but to exact bytes.

Credit here does not depend on agreement. Where this work reproduces a published
result it says so, and where it fails to reproduce one it says that too, under the
same names, because a refutation rests on the original as much as a confirmation does.

\begin{description}
\item[Registry commit] \texttt{2aae1103d28d9b919eae59cc9062aeaffe1fa0b2}
\item[Dated] 2026-10-01
\item[Inventory] \texttt{MANIFEST.tsv}, signed; the hashes below are its rows
\end{description}

A reader who disagrees with a number in this research paper can ask which copy it was checked
against, and the hash answers exactly. Two copies of a paper are rarely the same
bytes.
\par\medskip\raggedright}
\KOMAoptions{bibliography=totoc}
\begin{thebibliography}{6}

\bibitem{src:Langmuir-1919}
I. Langmuir. \emph{The arrangement of electrons in atoms and molecules}. Journal of the American Chemical Society 41(6) 868-934, \texttt{doi:10.1021/\allowbreak{}ja02227a002}. 1919.

\bibitem{src:Lewis-1916}
G. N. Lewis. \emph{The atom and the molecule}. Journal of the American Chemical Society 38(4) 762-785, \texttt{doi:10.1021/\allowbreak{}ja02261a002}. 1916.

\bibitem{src:Mendeleev-1869}
D. Mendeleev. 1869.

\bibitem{src:Moseley-1913}
H. Moseley. \emph{The high-frequency spectra of the elements}. The London, Edinburgh, and Dublin Philosophical Magazine and Journal of Science 26(156) 1024-1034, \texttt{doi:10.1080/\allowbreak{}14786441308635052}. 1913.

\bibitem{src:Pauli-1925}
W. Pauli. \emph{Über den Zusammenhang des Abschlusses der Elektronengruppen im Atom mit der Komplexstruktur der Spektren}. Zeitschrift für Physik 31, 765-783, \texttt{doi:10.1007/\allowbreak{}bf02980631}. 1925.

\bibitem{src:Pauling-1932}
L. Pauling. \emph{The nature of the chemical bond}. Journal of the American Chemical Society 54(9) 3570-3582, \texttt{doi:10.1021/\allowbreak{}ja01348a011}. 1932.

\end{thebibliography}
